\chapter{Boolean Analysis --- Circuit}
%%% generated by the blueprint pipeline from TCSlib.BooleanAnalysis.Switching.Circuit
\section{Overview}
This module sets up the basic combinatorial data of Boolean circuits used in the
switching lemma development: literals, general AND/OR circuit trees and their
size/depth/fanin measures, normal-form alternating circuits ($\mathsf{NAnd}$ and
$\mathsf{NOr}$) with their normalization maps and semantics/size/literal-count
preservation results, as well as the auxiliary types (literals, terms, DNF/CNF, and
decision trees) needed elsewhere in the switching-lemma proof.

\section{Declarations}

\begin{definition}[Evaluation of a literal]
\lean{BoolCircuit.Lit.eval}
\statementsource{boolean-ch04-dnf-switching-lmn}{pdf-fc447f326910-p001-b004}
\label{BoolCircuit.Lit.eval}
\leanok
\uses{BoolCircuit.Lit}
For a literal $l$ on $n$ variables and an assignment $x : \mathrm{Fin}\,n \to \bbf_2$,
the value $l.\mathrm{eval}\,x$ is $x_{l.\mathrm{idx}}$ when $l$ is positive
($l.\mathrm{sign} = \mathrm{true}$) and $\lnot x_{l.\mathrm{idx}}$ when $l$ is negated.
\end{definition}

\begin{definition}[Boolean circuit tree]
\lean{BoolCircuit.Circuit}
\statementsource{boolean-ch04-dnf-switching-lmn}{pdf-fc447f326910-p011-b004}
\label{BoolCircuit.Circuit}
\leanok
\uses{BoolCircuit.Lit}
A Boolean circuit on $n$ variables is a tree whose leaves are literals
(\texttt{lit}\,$l$) and whose internal nodes (\texttt{node}\,$isAnd$\,$children$) are
AND gates when $isAnd = \mathrm{true}$ and OR gates when $isAnd = \mathrm{false}$,
applied to a list of child circuits. No alternation or deduplication constraint is
imposed.
\end{definition}

\begin{theorem}[Structural induction on Boolean circuits]
\lean{BoolCircuit.Circuit.ind}
\label{BoolCircuit.Circuit.ind}
\leanok
\uses{BoolCircuit.Circuit, BoolCircuit.Circuit.eval}
\difficulty{5}
Let $P$ be a property of Boolean circuits on $n$ variables. Suppose that $P$ holds for
every literal circuit, and that whenever $P$ holds for every circuit in a list $cs$ of
children, it also holds for the gate circuit $\mathrm{node}(isAnd, cs)$ formed from $cs$
(an AND gate when $isAnd = \mathrm{true}$, an OR gate when $isAnd = \mathrm{false}$).
Then $P$ holds for every Boolean circuit on $n$ variables.
\end{theorem}

\begin{definition}[Evaluation of a circuit]
\lean{BoolCircuit.Circuit.eval}
\statementsource{boolean-ch04-dnf-switching-lmn}{pdf-fc447f326910-p011-b004}
\label{BoolCircuit.Circuit.eval}
\leanok
\uses{BoolCircuit.Circuit, BoolCircuit.Lit.eval}
Evaluates a circuit under an assignment $x$: a literal evaluates via
$\mathrm{Lit.eval}$, an AND node is the conjunction (\texttt{foldr} with \texttt{\&\&},
unit \texttt{true}) of its children's values, and an OR node is the disjunction
(\texttt{foldr} with \texttt{||}, unit \texttt{false}).
\end{definition}

\begin{definition}[Literal count of a circuit]
\lean{BoolCircuit.Circuit.litCount}
\label{BoolCircuit.Circuit.litCount}
\leanok
\uses{BoolCircuit.Circuit}
The number of literal occurrences in a circuit: a literal contributes $1$ and a node
contributes the sum of the literal counts of its children.
\end{definition}

\begin{definition}[Depth of a circuit]
\lean{BoolCircuit.Circuit.depth}
\statementsource{boolean-ch04-dnf-switching-lmn}{pdf-fc447f326910-p011-b004}
\label{BoolCircuit.Circuit.depth}
\leanok
\uses{BoolCircuit.Circuit}
The depth of a circuit (longest root-to-leaf path): a literal has depth $0$ and a node
has depth $1$ plus the maximum depth among its children.
\end{definition}

\begin{definition}[Size of a circuit]
\lean{BoolCircuit.Circuit.size}
\statementsource{boolean-ch04-dnf-switching-lmn}{pdf-fc447f326910-p012-b001}
\label{BoolCircuit.Circuit.size}
\leanok
\uses{BoolCircuit.Circuit}
The total number of nodes of a circuit (internal gates plus literal leaves): a literal
has size $1$ and a node has size $1$ plus the sum of the sizes of its children.
\end{definition}

\begin{definition}[Maximum depth over a list of circuits]
\lean{BoolCircuit.Circuit.maxDepth}
\label{BoolCircuit.Circuit.maxDepth}
\leanok
\uses{BoolCircuit.Circuit, BoolCircuit.Circuit.depth}
The maximum depth over a list $cs$ of circuits (with $0$ for the empty list), used to
compute the depth at a node.
\end{definition}

\begin{definition}[Sum of sizes over a list of circuits]
\lean{BoolCircuit.Circuit.sumSize}
\label{BoolCircuit.Circuit.sumSize}
\leanok
\uses{BoolCircuit.Circuit, BoolCircuit.Circuit.size}
The sum of the sizes over a list $cs$ of circuits, used to compute the size at a node.
\end{definition}

\begin{definition}[Maximum fanin of a circuit]
\lean{BoolCircuit.Circuit.maxFanin}
\statementsource{boolean-ch04-dnf-switching-lmn}{pdf-fc447f326910-p012-b001}
\label{BoolCircuit.Circuit.maxFanin}
\leanok
\uses{BoolCircuit.Circuit}
The maximum fanin of a circuit: a literal has fanin $0$, and a node's value is the
maximum of the number of its children and the maximal fanin occurring recursively among
its children.
\end{definition}

\begin{definition}[Normal-form AND circuit]
\lean{BoolCircuit.NAndCircuit}
\statementsource{boolean-ch04-dnf-switching-lmn}{pdf-fc447f326910-p011-b004}
\label{BoolCircuit.NAndCircuit}
\leanok
\uses{BoolCircuit.Lit}
A normal-form AND circuit on $n$ variables: either a base \texttt{clause} given by a list
of literals whose variable indices are pairwise distinct (a \texttt{Nodup} proof on
$lits.\mathrm{map}\;\mathrm{Lit.idx}$), or a \texttt{node} that ANDs together a list of
normal-form OR circuits. Together with $\mathrm{NOrCircuit}$ this forms a strictly
alternating AND/OR normal form.
\end{definition}

\begin{definition}[Evaluation of a normal-form AND circuit]
\lean{BoolCircuit.NAndCircuit.eval}
\statementsource{boolean-ch04-dnf-switching-lmn}{pdf-fc447f326910-p011-b004}
\label{BoolCircuit.NAndCircuit.eval}
\leanok
\uses{BoolCircuit.Lit.eval, BoolCircuit.NAndCircuit}
Evaluates a normal-form AND circuit under an assignment $x$: a clause is the conjunction
of its literals' values (unit \texttt{true}), and a node is the conjunction of its child
OR circuits' values.
\end{definition}

\begin{definition}[Evaluation of a normal-form OR circuit]
\lean{BoolCircuit.NOrCircuit.eval}
\statementsource{boolean-ch04-dnf-switching-lmn}{pdf-fc447f326910-p011-b004}
\label{BoolCircuit.NOrCircuit.eval}
\leanok
\uses{BoolCircuit.Lit.eval, BoolCircuit.NAndCircuit.eval}
Evaluates a normal-form OR circuit under an assignment $x$: a clause is the disjunction
of its literals' values (unit \texttt{false}), and a node is the disjunction of its child
AND circuits' values.
\end{definition}

\begin{definition}[Literal count of a normal-form AND circuit]
\lean{BoolCircuit.NAndCircuit.litCount}
\label{BoolCircuit.NAndCircuit.litCount}
\leanok
\uses{BoolCircuit.NAndCircuit}
The number of literal occurrences in a normal-form AND circuit: a clause contributes the
length of its literal list, and a node contributes the sum over its child OR circuits.
\end{definition}

\begin{definition}[Literal count of a normal-form OR circuit]
\lean{BoolCircuit.NOrCircuit.litCount}
\label{BoolCircuit.NOrCircuit.litCount}
\leanok
\uses{BoolCircuit.NAndCircuit.litCount}
The number of literal occurrences in a normal-form OR circuit: a clause contributes the
length of its literal list, and a node contributes the sum over its child AND circuits.
\end{definition}

\begin{definition}[Size of a normal-form AND circuit]
\lean{BoolCircuit.NAndCircuit.size}
\statementsource{boolean-ch04-dnf-switching-lmn}{pdf-fc447f326910-p012-b001}
\label{BoolCircuit.NAndCircuit.size}
\leanok
\uses{BoolCircuit.NAndCircuit}
The total node count of a normal-form AND circuit: a clause has size $1$, and a node has
size $1$ plus the sum of the sizes of its child OR circuits.
\end{definition}

\begin{definition}[Size of a normal-form OR circuit]
\lean{BoolCircuit.NOrCircuit.size}
\statementsource{boolean-ch04-dnf-switching-lmn}{pdf-fc447f326910-p012-b001}
\label{BoolCircuit.NOrCircuit.size}
\leanok
\uses{BoolCircuit.NAndCircuit.size}
The total node count of a normal-form OR circuit: a clause has size $1$, and a node has
size $1$ plus the sum of the sizes of its child AND circuits.
\end{definition}

\begin{definition}[Depth of a normal-form AND circuit]
\lean{BoolCircuit.NAndCircuit.depth}
\statementsource{boolean-ch04-dnf-switching-lmn}{pdf-fc447f326910-p011-b004}
\label{BoolCircuit.NAndCircuit.depth}
\leanok
\uses{BoolCircuit.NAndCircuit}
The depth of a normal-form AND circuit: a clause has depth $0$, and a node has depth $1$
plus the maximum depth among its child OR circuits.
\end{definition}

\begin{definition}[Depth of a normal-form OR circuit]
\lean{BoolCircuit.NOrCircuit.depth}
\statementsource{boolean-ch04-dnf-switching-lmn}{pdf-fc447f326910-p011-b004}
\label{BoolCircuit.NOrCircuit.depth}
\leanok
\uses{BoolCircuit.NAndCircuit.depth}
The depth of a normal-form OR circuit: a clause has depth $0$, and a node has depth $1$
plus the maximum depth among its child AND circuits.
\end{definition}

\begin{theorem}[Distinct variable indices in an AND clause]
\lean{BoolCircuit.NAndCircuit.clause_nodup}
\label{BoolCircuit.NAndCircuit.clause_nodup}
\leanok
\uses{BoolCircuit.Lit, BoolCircuit.NAndCircuit}
\difficulty{1}
Fix $n$, and let $lits$ be a list of literals over the variables $\{0,\dots,n-1\}$ whose
underlying variable indices, obtained by taking the index of each literal in $lits$, are
pairwise distinct. If a normal-form AND circuit on $n$ variables equals the base clause
built from $lits$ together with this distinctness witness, then the list of variable
indices of $lits$ has no repeated entries.
\end{theorem}

\begin{theorem}[Distinct variable indices in an OR clause]
\lean{BoolCircuit.NOrCircuit.clause_nodup}
\label{BoolCircuit.NOrCircuit.clause_nodup}
\leanok
\uses{BoolCircuit.Lit}
\difficulty{1}
Fix a natural number $n$, and let $\ell$ be a list of literals over $n$ variables, each
literal being a pair consisting of a variable index in $\mathrm{Fin}\,n$ and a sign in
$\mathrm{Bool}$. Suppose the variable indices occurring in $\ell$ are pairwise distinct,
and let $c$ be the OR-clause circuit built from $\ell$ under this distinctness
hypothesis. Then the list of variable indices of $\ell$ is duplicate-free, i.e. its
entries are pairwise distinct.
\end{theorem}

\begin{theorem}[Distinct indices determine literals in a list]
\lean{BoolCircuit.Lit.eq_of_idx_eq_of_mem_nodup}
\label{BoolCircuit.Lit.eq_of_idx_eq_of_mem_nodup}
\leanok
\uses{BoolCircuit.Lit}
\difficulty{4}
Let $\ell_1,\dots,\ell_m$ be a list of circuit literals over $n$ variables whose
variable indices are pairwise distinct. Then any two literals $l_1$ and $l_2$ in this
list that have the same variable index are equal; that is, if $l_1$ and $l_2$ both occur
in the list and $l_1.\mathrm{idx} = l_2.\mathrm{idx}$, then $l_1 = l_2$.
\end{theorem}

\begin{definition}[Normalize a circuit to AND-normal form]
\lean{BoolCircuit.Circuit.toNAnd}
\statementsource{boolean-ch04-dnf-switching-lmn}{pdf-fc447f326910-p012-b002}
\label{BoolCircuit.Circuit.toNAnd}
\leanok
\uses{BoolCircuit.Circuit, BoolCircuit.NAndCircuit}
Converts a general circuit into a normal-form AND circuit: a literal becomes a singleton
clause; an AND node maps its children through $\mathrm{toNOr}$; and an OR node is wrapped
as a single OR node built from the children's $\mathrm{toNAnd}$ images, enforcing
alternation.
\end{definition}

\begin{definition}[Normalize a circuit to OR-normal form]
\lean{BoolCircuit.Circuit.toNOr}
\statementsource{boolean-ch04-dnf-switching-lmn}{pdf-fc447f326910-p012-b002}
\label{BoolCircuit.Circuit.toNOr}
\leanok
\uses{BoolCircuit.Circuit, BoolCircuit.Circuit.toNAnd}
Converts a general circuit into a normal-form OR circuit: a literal becomes a singleton
clause; an OR node maps its children through $\mathrm{toNAnd}$; and an AND node is wrapped
as a single AND node built from the children's $\mathrm{toNOr}$ images, enforcing
alternation.
\end{definition}

\begin{theorem}[Mapped conjunction fold equals the original]
\lean{BoolCircuit.foldr_and_map}
\label{BoolCircuit.foldr_and_map}
\leanok
\difficulty{3}
Let $\alpha$ and $\beta$ be types, let
$f\colon\alpha\to\{\mathrm{true},\mathrm{false}\}$ and
$g\colon\beta\to\{\mathrm{true},\mathrm{false}\}$ be Boolean-valued functions, let
$h\colon\alpha\to\beta$, and let $cs$ be a finite list of elements of $\alpha$. Suppose
that $g(h(c)) = f(c)$ for every $c$ appearing in $cs$. Then the Boolean conjunction of
$g$ over the mapped list, $\bigwedge_{c\in cs} g(h(c))$, equals the Boolean conjunction
of $f$ over $cs$, namely $\bigwedge_{c\in cs} f(c)$, where in each case the conjunction
is folded from the right starting from $\mathrm{true}$.
\end{theorem}

\begin{theorem}[Disjunctive fold commutes with a compatible map]
\lean{BoolCircuit.foldr_or_map}
\label{BoolCircuit.foldr_or_map}
\leanok
\difficulty{3}
Let $f : \alpha \to \{\text{true},\text{false}\}$ and $g : \beta \to
\{\text{true},\text{false}\}$ be Boolean-valued functions, let $h : \alpha \to \beta$,
and let $c_1, \dots, c_n$ be a finite list of elements of $\alpha$. If $g(h(c)) = f(c)$
for every element $c$ of the list, then the disjunction $g(h(c_1)) \vee \cdots \vee
g(h(c_n))$ of the values of $g$ over the mapped list equals the disjunction $f(c_1) \vee
\cdots \vee f(c_n)$ of the values of $f$ over the original list, where each empty
disjunction is taken to be $\text{false}$.
\end{theorem}

\begin{theorem}[Sums agree under a compatible relabelling]
\lean{BoolCircuit.foldr_add_map}
\label{BoolCircuit.foldr_add_map}
\leanok
\difficulty{3}
Let $f\colon A \to \bbn$ and $g\colon B \to \bbn$ be functions into the natural numbers,
let $h\colon A \to B$ be a map, and let $c_1,\dots,c_m$ be a finite list of elements of
$A$. If $g(h(c)) = f(c)$ for every entry $c$ of the list, then
\[
\sum_{i=1}^{m} g\bigl(h(c_i)\bigr) \;=\; \sum_{i=1}^{m} f(c_i),
\]
where each sum is taken over the list in order with initial value $0$.
\end{theorem}

\begin{theorem}[Sums commute with a mapped pointwise bound]
\lean{BoolCircuit.foldr_add_map_le}
\label{BoolCircuit.foldr_add_map_le}
\leanok
\difficulty{3}
Let $f\colon\alpha\to\bbn$, $g\colon\beta\to\bbn$, and $h\colon\alpha\to\beta$ be
functions into the natural numbers, let $cs$ be a finite list of elements of $\alpha$,
and let $k\in\bbn$. If $g(h(c)) \le k\, f(c)$ holds for every entry $c$ of $cs$, then,
summing over the entries of $cs$ (with multiplicity),
\[ \sum_{c \in cs} g\bigl(h(c)\bigr) \;\le\; k \sum_{c \in cs} f(c). \]
\end{theorem}

\begin{theorem}[Normalization preserves circuit semantics]
\lean{BoolCircuit.toNAnd_toNOr_eval}
\label{BoolCircuit.toNAnd_toNOr_eval}
\leanok
\uses{BoolCircuit.Circuit, BoolCircuit.Circuit.eval, BoolCircuit.Circuit.ind, BoolCircuit.Circuit.toNAnd, BoolCircuit.Circuit.toNOr, BoolCircuit.NAndCircuit.eval, BoolCircuit.NOrCircuit.eval}
\difficulty{5}
Let $c$ be a Boolean circuit on $n$ variables and let $x : \mathrm{Fin}\,n \to \bbf_2$
be an assignment. Then both alternating normal forms of $c$ compute the same value as
$c$ itself: the normal-form AND circuit obtained by normalizing $c$ evaluates at $x$ to
$c$'s value, and likewise the normal-form OR circuit obtained by normalizing $c$
evaluates at $x$ to $c$'s value.
\end{theorem}

\begin{theorem}[Correctness of AND-normal-form conversion]
\lean{BoolCircuit.toNAnd_eval}
\label{BoolCircuit.toNAnd_eval}
\leanok
\uses{BoolCircuit.Circuit, BoolCircuit.Circuit.eval, BoolCircuit.Circuit.toNAnd, BoolCircuit.NAndCircuit.eval, BoolCircuit.toNAnd_toNOr_eval}
\difficulty{1}
Let $c$ be a Boolean circuit on $n$ variables — a tree of literals combined by AND and
OR gates — and let $x : \mathrm{Fin}\,n \to \bbf_2$ be an assignment to those variables.
Converting $c$ into strictly alternating AND-normal form and evaluating the resulting
normal-form AND circuit at $x$ yields the same Boolean value as evaluating $c$ directly
at $x$.
\end{theorem}

\begin{theorem}[OR-normalization preserves semantics]
\lean{BoolCircuit.toNOr_eval}
\label{BoolCircuit.toNOr_eval}
\leanok
\uses{BoolCircuit.Circuit, BoolCircuit.Circuit.eval, BoolCircuit.Circuit.toNOr, BoolCircuit.NOrCircuit.eval, BoolCircuit.toNAnd_toNOr_eval}
\difficulty{1}
Let $c$ be a Boolean circuit on $n$ variables and let $x : \mathrm{Fin}\,n \to \bbf_2$
be an assignment. Then the normal-form OR circuit obtained by OR-normalizing $c$
computes the same Boolean value as $c$ itself: evaluated at $x$, it agrees with the
value of $c$ at $x$.
\end{theorem}

\begin{theorem}[Normalization preserves the literal count]
\lean{BoolCircuit.toNAnd_toNOr_litCount}
\label{BoolCircuit.toNAnd_toNOr_litCount}
\leanok
\uses{BoolCircuit.Circuit, BoolCircuit.Circuit.ind, BoolCircuit.Circuit.litCount, BoolCircuit.Circuit.toNAnd, BoolCircuit.Circuit.toNOr, BoolCircuit.NAndCircuit.litCount, BoolCircuit.NOrCircuit.litCount}
\difficulty{5}
Let $c$ be a Boolean circuit on $n$ variables, whose literal count is the number of
literal occurrences it contains (a leaf literal counting once, and an AND or OR gate
counting the total over its children). Converting $c$ to strictly alternating AND/OR
normal form preserves this count: both the AND-normal form of $c$ and the OR-normal form
of $c$ have the same number of literal occurrences as $c$ itself.
\end{theorem}

\begin{theorem}[AND-normalization preserves the literal count]
\lean{BoolCircuit.toNAnd_litCount}
\label{BoolCircuit.toNAnd_litCount}
\leanok
\uses{BoolCircuit.Circuit, BoolCircuit.Circuit.litCount, BoolCircuit.Circuit.toNAnd, BoolCircuit.NAndCircuit.litCount, BoolCircuit.toNAnd_toNOr_litCount}
\difficulty{1}
Let $c$ be a Boolean circuit on $n$ variables, and let $c'$ be the normal-form AND
circuit obtained by normalizing it (a literal becomes a singleton clause, an AND gate
normalizes each child to OR-normal form, and an OR gate is wrapped as a single OR node
over the AND-normalizations of its children). Then $c'$ contains exactly as many literal
occurrences as $c$: the number of literals in $c'$ equals the number of literals in $c$.
\end{theorem}

\begin{theorem}[OR-normalization preserves literal count]
\lean{BoolCircuit.toNOr_litCount}
\label{BoolCircuit.toNOr_litCount}
\leanok
\uses{BoolCircuit.Circuit, BoolCircuit.Circuit.litCount, BoolCircuit.Circuit.toNOr, BoolCircuit.NOrCircuit.litCount, BoolCircuit.toNAnd_toNOr_litCount}
\difficulty{1}
For every Boolean circuit $c$ on $n$ variables, converting $c$ to OR-normal form leaves
its total number of literal occurrences unchanged: the literal count of the normal-form
OR circuit obtained from $c$ equals the literal count of $c$ itself.
\end{theorem}

\begin{theorem}[Normalization at most doubles circuit size]
\lean{BoolCircuit.toNAnd_toNOr_size_le}
\statementsource{boolean-ch04-dnf-switching-lmn}{pdf-fc447f326910-p012-b002}
\label{BoolCircuit.toNAnd_toNOr_size_le}
\leanok
\uses{BoolCircuit.Circuit, BoolCircuit.Circuit.ind, BoolCircuit.Circuit.size, BoolCircuit.Circuit.toNAnd, BoolCircuit.Circuit.toNOr, BoolCircuit.NAndCircuit.size, BoolCircuit.NOrCircuit.size, BoolCircuit.foldr_add_map_le}
\difficulty{6}
Let $c$ be a Boolean circuit on $n$ variables, and write $\mathrm{size}(c)$ for its
total node count. Converting $c$ into either alternating normal form increases its size
by at most a factor of two: the normal-form AND circuit obtained by normalizing $c$ has
size at most $2\,\mathrm{size}(c)$, and the normal-form OR circuit obtained by
normalizing $c$ likewise has size at most $2\,\mathrm{size}(c)$.
\end{theorem}

\begin{theorem}[AND-normalization at most doubles circuit size]
\lean{BoolCircuit.toNAnd_size_le}
\label{BoolCircuit.toNAnd_size_le}
\leanok
\uses{BoolCircuit.Circuit, BoolCircuit.Circuit.size, BoolCircuit.Circuit.toNAnd, BoolCircuit.NAndCircuit.size, BoolCircuit.toNAnd_toNOr_size_le}
\difficulty{1}
Let $c$ be a Boolean circuit on $n$ variables, where the size of a circuit is its total
number of nodes, counting internal AND/OR gates together with literal leaves.
Normalizing $c$ into strictly alternating AND-normal form — a literal becoming a
singleton clause, an AND gate normalizing each child into OR-normal form, and an OR gate
becoming a single wrapping OR node over its children's AND-normal forms — yields a
normal-form AND circuit whose size is at most twice the size of $c$.
\end{theorem}

\begin{theorem}[OR-normalization at most doubles size]
\lean{BoolCircuit.toNOr_size_le}
\statementsource{boolean-ch04-dnf-switching-lmn}{pdf-fc447f326910-p012-b002}
\label{BoolCircuit.toNOr_size_le}
\leanok
\uses{BoolCircuit.Circuit, BoolCircuit.Circuit.size, BoolCircuit.Circuit.toNOr, BoolCircuit.NOrCircuit.size, BoolCircuit.toNAnd_toNOr_size_le}
\difficulty{1}
Let $c$ be a Boolean circuit on $n$ variables, and let $\mathrm{size}(c)$ denote its
total node count (literal leaves plus internal AND/OR gates). Write $c'$ for the
OR-normal form of $c$, the normal-form OR circuit obtained by rewriting $c$ into
strictly alternating AND/OR normal form. Then the size of $c'$ is at most twice the size
of $c$:
\[
\mathrm{size}(c') \le 2\,\mathrm{size}(c).
\]
\end{theorem}

\begin{definition}[Forget AND-normal form back to a circuit]
\lean{BoolCircuit.NAndCircuit.toCircuit}
\label{BoolCircuit.NAndCircuit.toCircuit}
\leanok
\uses{BoolCircuit.Circuit, BoolCircuit.NAndCircuit}
The forgetful map turning a normal-form AND circuit into a general circuit: a clause
becomes an AND node over its literals, and a node becomes an AND node over the
circuit images of its child OR circuits.
\end{definition}

\begin{definition}[Forget OR-normal form back to a circuit]
\lean{BoolCircuit.NOrCircuit.toCircuit}
\label{BoolCircuit.NOrCircuit.toCircuit}
\leanok
\uses{BoolCircuit.Circuit, BoolCircuit.NAndCircuit.toCircuit}
The forgetful map turning a normal-form OR circuit into a general circuit: a clause
becomes an OR node over its literals, and a node becomes an OR node over the circuit
images of its child AND circuits.
\end{definition}

\begin{definition}[Single-variable AND circuit]
\lean{BoolCircuit.NAndCircuit.ofVar}
\label{BoolCircuit.NAndCircuit.ofVar}
\leanok
\uses{BoolCircuit.NAndCircuit}
The normal-form AND circuit consisting of the single positive literal on variable $i$.
\end{definition}

\begin{definition}[Single-variable OR circuit]
\lean{BoolCircuit.NOrCircuit.ofVar}
\label{BoolCircuit.NOrCircuit.ofVar}
\leanok
The normal-form OR circuit consisting of the single positive literal on variable $i$.
\end{definition}

\begin{definition}[Constant-true AND circuit]
\lean{BoolCircuit.NAndCircuit.constTrue}
\label{BoolCircuit.NAndCircuit.constTrue}
\leanok
\uses{BoolCircuit.NAndCircuit}
The constant-true AND circuit, given by the empty clause (an empty conjunction).
\end{definition}

\begin{definition}[Constant-false OR circuit]
\lean{BoolCircuit.NOrCircuit.constFalse}
\label{BoolCircuit.NOrCircuit.constFalse}
\leanok
The constant-false OR circuit, given by the empty clause (an empty disjunction).
\end{definition}

\begin{definition}[Evaluation of a literal]
\lean{Literal.eval}
\statementsource{boolean-ch04-dnf-switching-lmn}{pdf-fc447f326910-p001-b004}
\label{Literal.eval}
\leanok
\uses{Literal}
For a literal $l$ with variable $l.\mathrm{var}$ and polarity $l.\mathrm{neg}$ and an
assignment $x$, the value is $x_{l.\mathrm{var}}$ for a positive literal and
$\lnot x_{l.\mathrm{var}}$ for a negated one.
\end{definition}

\begin{definition}[Width of a CNF formula]
\lean{CNF.width}
\statementsource{boolean-ch04-dnf-switching-lmn}{
  pdf-fc447f326910-p002-b004,
  pdf-fc447f326910-p002-b002
}
\statementsource{switching-lemma}{pdf-b5e074215b9e-p001-b004}
\label{CNF.width}
\leanok
\uses{CNF, Term.width}
The width of a CNF formula $c$, defined as the maximum width over its clauses (and $0$
for the empty formula).
\end{definition}

\begin{definition}[Evaluation of a decision tree]
\lean{DecisionTree.eval}
\statementsource{switching-lemma}{pdf-b5e074215b9e-p001-b005}
\label{DecisionTree.eval}
\leanok
\uses{DecisionTree}
Evaluates a decision tree on an input $x$: a leaf returns its stored bit, and a branch
on variable $i$ follows the high subtree when $x_i$ is \texttt{true} and the low subtree
otherwise.
\end{definition}

\begin{definition}[Depth of a decision tree]
\lean{DecisionTree.depth}
\statementsource{switching-lemma}{pdf-b5e074215b9e-p001-b005}
\label{DecisionTree.depth}
\leanok
\uses{DecisionTree}
The depth of a decision tree (maximum path length to a leaf): a leaf has depth $0$ and a
branch has depth $1$ plus the maximum of its two subtrees' depths.
\end{definition}

\begin{definition}[Deepest path of a decision tree]
\lean{DecisionTree.deepPath}
\label{DecisionTree.deepPath}
\leanok
\uses{DecisionTree, DecisionTree.depth}
Extracts a deepest root-to-leaf path from a decision tree, at each branch descending into
the deeper subtree (ties broken toward the high subtree). It returns the list of
$(\text{queried variable}, \text{branch direction})$ pairs along that path.
\end{definition}

\begin{lemma}[Length of the deepest path equals the depth]
\lean{DecisionTree.length_deepPath}
\label{DecisionTree.length_deepPath}
\leanok
\uses{DecisionTree, DecisionTree.deepPath, DecisionTree.depth}
\difficulty{3}
Let $T$ be a decision tree on $n$ Boolean variables. Then the deepest root-to-leaf path
of $T$ has length equal to the depth of $T$; that is, the number of branch queries along
that path is exactly $\mathrm{depth}(T)$.
\end{lemma}

\begin{lemma}[Depth bound for the full decision tree]
\lean{buildFullDTree_depth}
\label{buildFullDTree_depth}
\leanok
\uses{DecisionTree.depth, buildFullDTree}
\difficulty{4}
Let $f\colon\{0,1\}^n \to \{0,1\}$ be a Boolean function on $n$ variables, let $k \le
n$, and let $acc$ be a partial assignment of the variables. Then the full decision tree
for $f$ that queries the variables $k, k+1, \dots, n-1$ in order starting from $acc$ has
depth at most $n - k$.
\end{lemma}

\begin{lemma}[Correctness of the full decision tree]
\lean{buildFullDTree_eval}
\label{buildFullDTree_eval}
\leanok
\uses{DecisionTree.eval, buildFullDTree}
\difficulty{5}
Let $f \colon (\bbf_2)^n \to \{\mathtt{true},\mathtt{false}\}$ be a Boolean function on
$n$ variables, let $k \le n$, and let $acc, x \colon \{0,\dots,n-1\} \to
\{\mathtt{true},\mathtt{false}\}$ be two assignments that agree on every coordinate
below $k$, that is $acc_i = x_i$ for all $i < k$. Then the full decision tree for $f$
that branches on variables $k, k+1, \dots, n-1$ starting from the accumulator $acc$
evaluates on input $x$ to $f(x)$.
\end{lemma}

%%%%%%%%%%%%%%%%%%%%%%%%%%%%%%%%%%%%%%%%%%%%%%%%%%%%%%%%%%%%%%%%%%%%%%%%%%%%%%%
\section{Base literal, formula, and decision-tree definitions}

\begin{definition}[Circuit literal]
\lean{BoolCircuit.Lit}
\statementsource{boolean-ch04-dnf-switching-lmn}{pdf-fc447f326910-p001-b004}
\label{BoolCircuit.Lit}
\leanok
A literal for Boolean circuits: a variable index $\mathtt{idx} : \mathrm{Fin}\,n$
together with a polarity $\mathtt{sign} : \mathrm{Bool}$ (true = positive occurrence).
\end{definition}

\begin{definition}[Literal]
\lean{Literal}
\statementsource{boolean-ch04-dnf-switching-lmn}{pdf-fc447f326910-p001-b004}
\label{Literal}
\leanok
A literal for switching-lemma formulas: a variable $\mathtt{var} : \mathrm{Fin}\,n$
with a negation flag $\mathtt{neg} : \mathrm{Bool}$ (true = negated literal).
\end{definition}

\begin{definition}[Term]
\lean{Term}
\statementsource{boolean-ch04-dnf-switching-lmn}{pdf-fc447f326910-p001-b004}
\statementsource{switching-lemma}{pdf-b5e074215b9e-p001-b004}
\label{Term}
\leanok
\uses{Literal}
A term is a conjunction of literals, represented as a list:
$\mathtt{Term}\,n = \mathtt{List}\,(\mathtt{Literal}\,n)$.
\end{definition}

\begin{definition}[Term width]
\lean{Term.width}
\statementsource{boolean-ch04-dnf-switching-lmn}{pdf-fc447f326910-p001-b004}
\statementsource{switching-lemma}{pdf-b5e074215b9e-p001-b004}
\label{Term.width}
\leanok
\uses{Term}
The width of a term is its number of literals (the list length).
\end{definition}

\begin{definition}[Term evaluation]
\lean{Term.eval}
\statementsource{boolean-ch04-dnf-switching-lmn}{pdf-fc447f326910-p001-b004}
\statementsource{switching-lemma}{pdf-b5e074215b9e-p001-b004}
\label{Term.eval}
\leanok
\uses{Term}
A term evaluates to true on input $x$ iff every literal in it holds under $x$
(conjunction semantics).
\end{definition}

\begin{definition}[DNF formula]
\lean{DNF}
\statementsource{boolean-ch04-dnf-switching-lmn}{pdf-fc447f326910-p001-b004}
\statementsource{switching-lemma}{pdf-b5e074215b9e-p001-b004}
\label{DNF}
\leanok
\uses{Term}
A DNF formula is a disjunction of terms, represented as a list:
$\mathtt{DNF}\,n = \mathtt{List}\,(\mathtt{Term}\,n)$.
\end{definition}

\begin{definition}[DNF width]
\lean{DNF.width}
\statementsource{boolean-ch04-dnf-switching-lmn}{pdf-fc447f326910-p002-b002}
\statementsource{switching-lemma}{pdf-b5e074215b9e-p001-b004}
\label{DNF.width}
\leanok
\uses{DNF, Term.width}
The width of a DNF formula is the maximum width of its terms ($0$ for the empty
formula).
\end{definition}

\begin{definition}[DNF evaluation]
\lean{DNF.eval}
\statementsource{boolean-ch04-dnf-switching-lmn}{pdf-fc447f326910-p001-b004}
\statementsource{switching-lemma}{pdf-b5e074215b9e-p001-b004}
\label{DNF.eval}
\leanok
\uses{DNF, Term.eval}
A DNF formula evaluates to true on input $x$ iff at least one of its terms holds
under $x$ (disjunction of conjunctions).
\end{definition}

\begin{definition}[CNF formula]
\lean{CNF}
\statementsource{boolean-ch04-dnf-switching-lmn}{pdf-fc447f326910-p002-b004}
\statementsource{switching-lemma}{pdf-b5e074215b9e-p001-b004}
\label{CNF}
\leanok
\uses{Term}
A CNF formula is a conjunction of clauses, each clause a disjunction of literals;
it is represented as $\mathtt{CNF}\,n = \mathtt{List}\,(\mathtt{Term}\,n)$ with
dual evaluation semantics.
\end{definition}

\begin{definition}[Clause evaluation]
\lean{CNF.evalClause}
\statementsource{boolean-ch04-dnf-switching-lmn}{pdf-fc447f326910-p002-b004}
\statementsource{switching-lemma}{pdf-b5e074215b9e-p001-b004}
\label{CNF.evalClause}
\leanok
\uses{CNF, Term}
A single CNF clause evaluates to true on input $x$ iff some literal in it holds
under $x$ (disjunction semantics).
\end{definition}

\begin{definition}[CNF evaluation]
\lean{CNF.eval}
\statementsource{boolean-ch04-dnf-switching-lmn}{pdf-fc447f326910-p002-b004}
\statementsource{switching-lemma}{pdf-b5e074215b9e-p001-b004}
\label{CNF.eval}
\leanok
\uses{CNF, CNF.evalClause}
A CNF formula evaluates to true on input $x$ iff every clause holds under $x$.
\end{definition}

\begin{definition}[Decision tree]
\lean{DecisionTree}
\statementsource{switching-lemma}{pdf-b5e074215b9e-p001-b005}
\label{DecisionTree}
\leanok
A decision tree on $n$ Boolean variables: a leaf $\mathtt{leaf}\,b$ outputs $b$,
and a branch $\mathtt{branch}\,i\,lo\,hi$ queries variable $i$, following $lo$ on
false and $hi$ on true.
\end{definition}

\begin{definition}[Full decision tree]
\lean{buildFullDTree}
\label{buildFullDTree}
\leanok
\uses{DecisionTree}
The complete decision tree for a function $f$ that queries variables
$k, k+1, \dots, n-1$ in order, starting from accumulator assignment $acc$, and
outputs $f$ applied to the completed assignment at each leaf.
\end{definition}

\begin{definition}[Decision-tree depth of a function]
\lean{dtDepth}
\statementsource{switching-lemma}{pdf-b5e074215b9e-p001-b005}
\label{dtDepth}
\leanok
\uses{DecisionTree, buildFullDTree}
The decision-tree depth of $f : (\mathrm{Fin}\,n \to \mathrm{Bool}) \to
\mathrm{Bool}$ is the least $d$ such that some decision tree of depth at most $d$
computes $f$; it is well defined since the full tree computes $f$ at depth $n$.
\end{definition}
